\subsection{Photometric perturbation}

We simulate test-time lighting changes by applying two single-parameter changes
to the validation images. Exposure and white-balance shifts are common camera
errors, and previous work has used both to test predictions under changed
capture conditions \cite{hendrycks2019benchmarking, afifi2019whatelse}; exposure
errors in particular are usually simulated in steps of exposure value (stops)
\cite{afifi2021exposure}. We apply both changes in linear light after undoing
the sRGB gamma curve, then convert the result back to sRGB, following the
standard practice of inverting the camera's processing before changing gains
\cite{brooks2019unprocessing}. A gain applied this way corresponds to what the
camera itself would have done, whereas the same gain applied to gamma-encoded
pixel values would not.
The exposure intervention multiplies every channel by $2^{k}$ for $k \in [-3,
3]$ stops. The white-balance intervention applies opposite gains to the red and
blue channels, $1 \pm 0.35w$ for $w \in [-1.5, 1.5]$, while keeping green fixed.
Positive $w$ produces a warm cast and negative $w$ a cool cast. Section~\ref{sec:colour}
measures ripeness along a green-to-red axis. A colour cast therefore moves the
apparent colour of every fruit along the axis used to assign its stage before
the detector processes the image.

The intervention recolours the pixels without changing the ground-truth boxes.
Exposure changes confidence scores across the predictions. We therefore run
each level once with a minimum confidence of 0.05 and score the resulting
predictions at every threshold in the sweep. The tables report a fixed
threshold of 0.25 alongside the best threshold for that level. A comparison at
only the fixed threshold would combine changes in calibration with changes in
robustness, as shown across datasets in Section~\ref{sec:calibration}.

\subsection{Training-time colour augmentation}

During training, Ultralytics YOLO changes the colour of every training image by
a random amount in HSV space, which represents hue, saturation, and value. This
form of photometric augmentation was introduced for the YOLO family in YOLOv4
\cite{bochkovskiy2020yolov4}. Three parameters set the maximum change. Their default values
are $h=0.015$, $s=0.7$ and $v=0.4$ \cite{ultralytics2026cfg}. This is the only colour
augmentation in the pipeline. We compare four configurations while keeping
everything else unchanged. The \texttt{stock} configuration uses the default
recipe. The \texttt{hsv\_v08} configuration doubles the brightness jitter to
$v=0.8$. The \texttt{hsv\_wide} configuration uses $v=0.8$ and widens hue to
$h=0.05$, which is the HSV change most similar to a colour cast. The
\texttt{no\_colour} configuration sets all three parameters to zero and provides
the comparison with no colour augmentation.

The test-time perturbation uses a different photometric process from the HSV
jitter used during training. An exposure change applied in linear light is
generated differently from the training-time noise. The evaluation therefore
tests whether the effect of HSV augmentation transfers to that distinct
process.
